\begin{table}[t]
\centering
\caption{Standardized Effect Sizes}
\label{tab:sde}
\footnotesize
\begin{tabular}{lcccccc}
\hline\hline
Outcome & $\hat{\beta}$ & SE & SD($Y$) & SDE & SE(SDE) & Class. \\
\hline
\textit{Panel A: Pooled} \\
Water temperature ($^{\circ}$C) & -0.242 & 0.362 & 3.01 & -0.080 & 0.120 & Mod.\ neg. \\
Dissolved oxygen (mg/L) & 0.214 & 0.152 & 1.19 & 0.180 & 0.128 & Large pos. \\
[0.5em]
\textit{Panel B: Heterogeneous (by proximity)} \\
Temperature, $\leq$ 10 km & 0.145 & 0.131 & 3.01 & 0.048 & 0.044 & Small pos. \\
Temperature, $>$ 10 km & -0.125 & 0.231 & 3.01 & -0.041 & 0.077 & Small neg. \\
\hline\hline
\end{tabular}
\par\vspace{0.5em}
\begin{minipage}{\textwidth}
\small
\textit{Notes:} \textbf{Country:} United States. \textbf{Research question:} Does removing a dam improve downstream water quality as measured by continuous sensor readings at USGS stream gauges? \textbf{Policy mechanism:} Dam removal physically eliminates a barrier that creates a warm-water impoundment, restoring natural flow regimes, thermal cycling, and sediment transport. \textbf{Outcome definition:} Panel A reports annual mean water temperature (degrees Celsius, USGS parameter 00010) and annual mean dissolved oxygen concentration (mg/L, parameter 00300) at stream gauges. Panel B splits the temperature outcome by gauge proximity to the removal site. \textbf{Treatment:} Binary---dam removed (year of removal from American Rivers Database). \textbf{Data:} American Rivers Dam Removal Database (1,341 removals, 2000--2020) matched to USGS NWIS daily stream gauge readings; 9,204 gauge-year observations for temperature, 2,719 for dissolved oxygen, 1995--2023. \textbf{Method:} Sun-Abraham (2021) interaction-weighted estimator with gauge and year fixed effects; standard errors clustered at the gauge level; overall ATT averages post-treatment event-time coefficients. \textbf{Sample:} USGS stream gauges within 20 km of a dam removal site (treated) or on rivers with no removals in the same states (control); restricted to gauges with at least 5 years of data and 6 months per year. SDE $= \hat{\beta} / \text{SD}(Y)$ where SD($Y$) is the pre-treatment standard deviation. Classification refers to magnitude, not statistical significance: Large ($|$SDE$| > 0.15$), Moderate ($0.05$--$0.15$), Small ($0.005$--$0.05$), Null ($< 0.005$).
\end{minipage}
\end{table}

